\documentclass[11pt]{article}

\usepackage[margin=25mm]{geometry}

\usepackage{amsmath,amssymb,fontspec,unicode-math}

\IfFontExistsTF{DejaVu Serif}{\setmainfont{DejaVu Serif}}{\setmainfont{Latin Modern Roman}}

\IfFontExistsTF{Latin Modern Math}{\setmathfont{Latin Modern Math}}{\setmathfont{DejaVu Math TeX Gyre}}

\usepackage{xurl,hyperref}

\setlength{\parindent}{0pt}\setlength{\parskip}{6pt}

\begin{document}

\title{Curriculum Design for Trajectory-Constrained Agent: Compressing Chain-of-Thought Tokens in LLMs}

\author{ELVIA preliminary AI analysis}\date{21 October 2026}\maketitle

Preliminary assessment: Major.\par

Independent scientific review remains pending. Author adoption has not been requested.

The curriculum theorem’s partition and value-accuracy reasoning fail on allowed discounts and policies. Its concentration proof also lacks a defined learner. The current zero-success code guard is a passing control.

\section*{Review scope}

Latest-version, claim-focused review of the cited author-paper passages and applicable current author-release paths. The stated counterexamples and scope controls were checked analytically; no author program, old verifier, training run or benchmark was executed.

\section*{Problem 1: Discounting invalidates the theorem’s first curriculum partition}

Source: Theorem 1, PDF p.6; Eq. (2), p. 5; Appendix C, current arXiv v1 PDF pp.17–18.

Assessment: supported; impact: major

\subsection*{Source claim (paraphrase)}

Theorem 1 permits \(\gamma\geq(1/2)^{1/H}\) and starts with threshold \(\beta_1=1/2\), while its proof gives successful depth-\(H\) trajectories undiscounted return one.

\subsection*{Analysis}

At \(H=2\), \(\gamma=\sqrt{1/2}\), every terminal success has discounted return \(1/2\). Accepting half the uniformly sampled leaves gives value \(1/4\), below \(\beta_1\); the teacher must retain all leaves. The next threshold \(0.45\) still requires all four. Thus the partition used to obtain the advertised curriculum does not occur.

\subsection*{Finding and scope}

The theorem’s construction fails for an included discount. A particular learner’s failure is not needed for this contradiction.

\subsection*{Proposed correction}

Scale curriculum thresholds by \(\gamma^H\), or explicitly normalize returns. This fixes the first partition arithmetic but does not supply the missing learning proof.

An explicit local correction is described; author adoption and a corresponding published erratum are not established.

\section*{Problem 2: Nodewise accuracy does not imply the stated value accuracy}

Source: Value objective, p. 4; Theorem 1 and Appendix C.

Assessment: supported; impact: major

\subsection*{Source claim (paraphrase)}

Theorem 1 equates final epsilon-suboptimality with every optimal-path action having probability at least \(1-\epsilon\).

\subsection*{Analysis}

The equality policy has terminal-success probability \((1-\epsilon)^H\), so its discounted gap is \(\gamma^H[1-(1-\epsilon)^H]\). At \(H=12\), \(\gamma=0.99\), \(\epsilon=2/13\), this is approximately \(0.76698\), exceeding \(0.15385\). The policy meets the stated nodewise requirement.

\subsection*{Finding and scope}

The claimed implication to the earlier value criterion is false.

\subsection*{Proposed correction}

To guarantee value gap at most \(\epsilon\), use a suitable per-depth error budget controlling \(1-\prod_h(1-\epsilon_h)\), rather than assigning the same target to every node.

An explicit local correction is described; author adoption and a corresponding published erratum are not established.

\section*{Problem 3: The concentration step does not specify a concentrating learner}

Source: Generic update in Algorithm 1; Theorem 1 proof sketch; Appendix C, current arXiv v1 PDF pp.17–18.

Assessment: supported; impact: minor

\subsection*{Source claim (paraphrase)}

The proof applies Hoeffding directly to updated policy probabilities, without defining the estimator and learning map that concentrate around an optimal probability.

\subsection*{Analysis}

An empirical reward average and a learned policy parameter are different random objects. Visit counts and the update rule need a linking argument. Even valid per-stage failure probability does not give the same joint failure probability: twelve independent successes of probability \(0.95\) have joint probability about \(0.54036\). Confidence allocation repairs only that arithmetic.

\subsection*{Finding and scope}

The advertised sample-complexity proof remains incomplete without a learner-specific argument.

\subsection*{Proposed correction}

If valid stagewise bounds are first proved, allocate failures with \(\sum_t\delta_t\leq\delta\), for example \(\delta_t=\delta/H\). Define the sampling and update mechanism before applying concentration.

An explicit local correction is described; author adoption and a corresponding published erratum are not established.

\section*{Problem 4: The zero-success formal rule differs from the current guard}

Source: Section 3.1, p. 5, Eq. (2); current constrained trainer find\_min\_alpha.

Assessment: supported; impact: minor

\subsection*{Source claim (paraphrase)}

With all rewards zero, Eq.(2)’s closest-target minimizer is \(\alpha^*\), while the prose and current implementation choose the upper budget \(U\).

\subsection*{Analysis}

Every budget is feasible when performance is zero, so the optimization alone selects the point closest to the target. The source explicitly checks `max\_perf == 0` and returns `upper\_limit`; the practical guard is already present.

\subsection*{Finding and scope}

The formal boundary needs alignment, but current training is not shown to select the wrong zero-success budget.

\subsection*{Proposed correction}

State the zero-performance branch piecewise before the closest-target optimization.

The current author code already contains the zero-success upper-budget guard; no missing-code correction is alleged.

\section*{Limits}

Preliminary AI assessment; independent human scientific adjudication remains pending.\par

An analytic counterexample or a release discrepancy does not by itself remeasure the paper’s reported experiments.\par

\section*{Source identity}

arXiv:2511.02690v1 (current author version; updated 2025-11-04T16:14:56Z)

\url{https://arxiv.org/pdf/2511.02690v1}

Reviewed PDF SHA-256: \nolinkurl{17d9c4884a929ce69e996a884994b6fc1b48b0b0cacfcfa0fd24a31244fc2216}

The selected author/publisher source is the same work as the original case, matched by exact title, author identity and official version channel. The exact selected PDF SHA-256 was recomputed in the latest-version pass for Issue \#27. The PDF is a retained author/publisher copy; its original retrieval time is stated only when supported by matching acquisition metadata. File modification time was not substituted.

\end{document}
